% ============================================================
% 05-conclusion.tex
% ============================================================
\section{Conclusion and Future Work}
\label{sec:conclusion}

This study investigated how centralized supervision and representation-specific evidence paths can support heterogeneous university counseling. The routing and robustness evaluations showed that the supervisor assigned queries to domain specialists with 97.0\% agent accuracy while maintaining bounded execution in the evaluated cases. In retrieval, the full configuration achieved the best Hit@1 ($0.9600$) and MRR@10 ($0.9700$), whereas BM25 retained the highest P@5 and Recall@5. Domain-stratified results further showed that the additional Hit@1 benefit of the full configuration was concentrated in Financial queries rather than distributed uniformly across domains. The reranking comparison improved early-rank ordering, while the E4-to-E5 comparison reflected the performance of the integrated routing, graph-grounding, and parent--child mapping bundle rather than an isolated routing effect. End-to-end ablations were also associated with lower selected context- and answer-quality metrics when graph grounding, reranking, or governance was removed.

Taken together, these findings suggest that no single representation or retrieval strategy dominates all university-counseling tasks. Relational curricula, structured financial information, and narrative regulations impose different evidence and reasoning requirements and are therefore better handled through separate but coordinated processing paths. In practice, the proposed architecture manages these paths within a unified deployment while restricting each specialist to its relevant knowledge sources and tools. The results support the integrated architecture in the evaluated setting, but do not establish that any individual component is universally optimal.

The evaluation remains limited to Can Tho University data, automated judging, offline knowledge schemas, small and imbalanced domain strata, and isolated robustness tests. Future work will evaluate transferability across institutions, conduct user studies with students and advisors, perform controlled computational-efficiency evaluation and model compression, employ cross-model evaluation to reduce judging bias, and compare \system{} with broader educational AI systems. Future work will also expand the production tool registry and evaluate matched single- and multi-agent configurations across increasing tool-set sizes, including production tools rather than synthetic distractors. Within the scope of this study, the central takeaway is not that one retrieval method should replace all others, but that institutional counseling is better supported when evidence acquisition and execution constraints are aligned with the structure of each domain.
